\section{Results}
\label{sec:results}

\subsection{Baselines vs.\ rule-conditioned NCA}
\label{sec:results-baselines}

Table~\ref{tab:baselines-microban} reports the head-to-head comparison
on \textsc{microban} (10 authored levels), trained for 10k updates per
seed at $d_h = 256$ with the canonical recipe (\S\ref{sec:experiments}).
We measure best train cross-entropy loss, end-of-rollout cell-error
rate under three policies (random, BFS-optimal, A*-optimal), and final
parameter count. The table reports a single seed; the multi-seed
synthetic-data comparison is reported in
Table~\ref{tab:baselines-synth}.

\begin{table}[t]
  \centering
  \caption{\textbf{Baselines vs.\ rule-conditioned NCA on
    \textsc{microban} (authored levels).} Single seed at 10k updates.
    Lower is better for loss and cell-error; A* first-divergence step
    is higher-is-better. Param count is in millions.}
  \label{tab:baselines-microban}
  \small
  \begin{tabular}{l c c c c c}
    \toprule
    Architecture & params (M) & best loss & BFS final-err & rand mean-err & A* first-div \\
    \midrule
    \textbf{NCA (shared, $T{=}4, n_r{=}T$)} & \textbf{1.94} & 0.106 & 0.34\% & 0.04\% & 64 \\
    NCA (per-step, $T{=}4, n_r{=}1$)        & 7.36          & 0.107 & 0.00\% & 0.09\% & none \\
    CNN (4-block ResNet)                    & 5.00          & 0.107 & 0.00\% & 0.09\% & none \\
    U-Net (2 levels)                        & 6.07          & 0.107 & 6.27\% & 3.00\% & 12 \\
    ViT (4 layers)                          & 3.99          & 0.107 & 11.67\% & 6.07\% & 1 \\
    \bottomrule
  \end{tabular}
\end{table}

Three findings on this controlled, memorizable single-game setting:
\textbf{(i)} the NCA with shared weights matches CNN / per-step NCA
on oracle-policy rollout error at $2.5{-}3.8\times$ fewer parameters,
\textbf{(ii)} U-Net and ViT both actively underperform despite
parameter counts comparable to or exceeding NCA-shared --- U-Net's
2-level downsampling loses fine-grained player position on small
grids, and ViT's global self-attention diverges within the first
rollout step (mean first-divergence step $0.8$), suggesting
locality plus iteration is a stronger inductive bias for grid-rule
dynamics than unrestricted global attention; and \textbf{(iii)} best
train loss is essentially identical (0.106--0.107) across all five
architectures, illustrating finding F3 (\S\ref{sec:results-F3}):
train loss is invisible to a $\geq 30\times$ swing in autoregressive
rollout error (0.34\% vs.\ 11.67\%).

\medskip\noindent
The full comparison on the multi-grid synthetic data
regime---where memorization is unavailable and rule learning is
required---is reported in Table~\ref{tab:baselines-synth} and
analyzed in \S\ref{sec:results-ceiling}.

\begin{table}[t]
  \centering
  \caption{\textbf{Baselines vs.\ rule-conditioned NCA on
    \textsc{sokoban\_basic} (multi-grid synthetic).} Single seed at 10k
    updates; widths $\{5{\times}5, 6{\times}6, 7{\times}7, 8{\times}8\}$
    with 256 synthetic levels per width. Memorization is unavailable;
    the model has to learn the rule. Architectures sorted by BFS final
    cell-error (lower is better). Multi-seed numbers are forthcoming.}
  \label{tab:baselines-synth}
  \small
  \begin{tabular}{l c c c c}
    \toprule
    Architecture & params (M) & best loss & BFS final-err & A* final-err \\
    \midrule
    \textbf{NCA (shared, $T{=}4, n_r{=}T$)} & \textbf{3.07} & \textbf{0.067} & \textbf{3.57\%} & \textbf{3.57\%} \\
    NCA (per-step, $T{=}4, n_r{=}1$)        & 8.49          & 0.068          & 4.76\%          & 4.76\% \\
    CNN (4-block ResNet)                    & 5.00          & 0.067          & 8.33\%          & 8.33\% \\
    U-Net (2 levels)                        & 6.07          & 0.070          & 9.52\%          & 9.52\% \\
    ViT (4 layers)                          & 3.99          & 0.085          & 11.90\%         & 11.90\% \\
    \bottomrule
  \end{tabular}
\end{table}

The synthetic-multi-grid setting is the headline comparison: per-state
memorization is impossible (256 levels per width $\times$ 4 widths
$=$ 1024 distinct levels share one rule set), so the model has to
\emph{learn the slide rule}. Three observations: \textbf{(i)} the
NCA with shared weights wins on BFS rollout error (3.57\%) at the
fewest parameters (3.07M), beating the same-depth feed-forward CNN
by $2.3\times$ in error and at $1.6\times$ fewer params;
\textbf{(ii)} U-Net (9.52\%) and ViT (11.90\%) both underperform
here as well, with U-Net's hierarchical bias actively losing in this
discrete-grid regime and ViT's global attention failing to even
\emph{fit} the train loss (0.085 vs.\ 0.067 for the other four);
\textbf{(iii)} train loss now does begin to discriminate
architectures (ViT's 0.085 outlier), but NCA-shared, NCA-per-step,
and CNN are again indistinguishable on train loss (0.067--0.068)
despite a $2.3\times$ difference in BFS rollout error --- F3 still
applies.

\paragraph{(F1) Pool features are load-bearing.}
A bare $3{\times}3$ convolution can only propagate one cell of context
per NCA step. Several PuzzleScript rule patterns reference cells
arbitrarily far apart along an axis or anywhere on the level
(\texttt{[X | ... | Y]}, \texttt{[X][Y]}). Without pool features the
NCA needs $\Omega(\max(H,W))$ steps just to surface non-local
information; with axis-cummax + global pool one step suffices.
Across the 19-game \textsc{scaling\_large} preset and a controlled
\textsc{collapse} no-pool ablation, turning pool off costs both
convergence speed and final loss, with the no-pool model collapsing
into the identity-prediction minimum even at $4{\times}$ the depth.

\paragraph{(F2) Shared NCA-body weights are the right inductive bias.}
The PuzzleScript engine has two natural levels of iteration: an
inner ordered list of rules applied once each per tick, and an outer
\texttt{again} loop that re-runs the rule list until the state stops
changing. We factor $T = n_\ell \cdot n_r$ accordingly. On a single-game
depth sweep ($T \in \{2, 4, 8, 16\}$), shared $T{=}8$ reaches
$5.5{\times}$ lower train loss with $7{\times}$ fewer parameters than
the per-step body at the same depth; shared $T{=}2$ at 14$\times$
fewer parameters than the deepest per-step variant matches it on
oracle-action rollouts. Hierarchical $n_r > 1$ sits between these
extremes and is the subject of an ongoing ablation
(\S\ref{sec:results-hierarchy}).

\paragraph{(F3) Train loss decouples from autoregressive rollout error.}
\label{sec:results-F3}
A model with the lowest train (per-step) loss is not the model with
the lowest autoregressive rollout error. Concretely, in the
over-capacity regime we observe per-step models with
$5{\times}$ lower train loss but $3{-}4{\times}$ higher 50-step
rollout error. The mechanism is small but systematic per-step errors
compounding destructively under autoregression. Reporting only
best-loss hides the rollout-drift cost of going deeper or wider.

\subsection{Pool-feature ablation}
\label{sec:results-pool}

[Figure: pool-on vs.\ pool-off on \textsc{collapse}, both at depths
$T \in \{2, 4, 8, 16\}$, reporting train loss + autoregressive
rollout error.]

\subsection{Hierarchical weight sharing}
\label{sec:results-hierarchy}

[Figure: $T{=}16$ varying $n_r \in \{1, 2, 4, 8, 16\}$ on
\textsc{microban} multi-grid synth and on the multi-game
\textsc{scaling\_14} preset. Reports change-error and parameter count.]

\subsection{Adaptive halting}
\label{sec:results-halting}

We compare four halting modes and a fixed-depth baseline on a custom
$1$-rule \textsc{varislide} game in which a single \texttt{again}-driven
slide rule pushes the player up to $d \in \{1, 2, 3, 4, 6, 8, 12, 16\}$
cells per action across grid widths $\{6, 8, 10, 12, 16\}$. Modes:
\emph{ponder} (PonderNet ponder loss with learned halt head + KL
prior), \emph{uniform} (per-step loss equally weighted across all
$T$ steps), \emph{argmax-ST} (straight-through estimator on
$\arg\max_k p_k$), and \emph{convergence-ST} (halt at inference when
consecutive predictions stop changing, no learned head).

[Figure: per-step error trajectory under each mode + the per-instance
halt distribution as a function of slide distance.]

Across these modes we observe two phenomena. First, on tasks where
the dynamics fit cleanly in $\sim$1--2 NCA steps (small-grid varislide,
\textsc{collapse}), all halting modes converge near the dynamics-floor
within their effective depth and drift past it; convergence-ST
recovers near-oracle performance via inference-time stopping.
Second, on the multi-grid synthetic regime where memorization is
unavailable, \emph{none} of the halting modes reliably learn the
underlying rule (\S\ref{sec:results-ceiling}); the halt head at best
degenerates to halt-at-$k{=}1$ + identity prediction.

\subsection{Multi-grid synthetic data ceiling}
\label{sec:results-ceiling}

When we move from a single authored level to the multi-grid synthetic
data regime ($\{6,8,10,12,16\}$ widths $\times$ 64 levels each),
train change-error climbs by $1000{\times}$ ($5 \times 10^{-5}
\to 0.05$--$0.10$), confirming that the model can no longer
1-shot-memorize per-state outcomes. The model attempts to learn the
rule but plateaus: 50k updates at $d_h{=}256$ (5$\times$ more updates,
2$\times$ wider) does not move the needle. Multi-seed verification at
$T \in \{8, 16, 32, 64\}$ shows per-seed argmax variance dwarfs the
recipe signal and best-loss is essentially identical across all
configurations. We discuss candidate fixes in \S\ref{sec:discussion}.

\subsection{Multi-game scaling}
\label{sec:results-scaling}

[Per-game change-error on \textsc{scaling\_14} for our model and the
three baselines; held-out evaluation on a few games not seen during
training.]
